% ============================================================================
% CHAPTER 4: Stempel-Wanderung
% ============================================================================
\renewcommand{\chaptertag}{Bildfilter}
\chapter{Stempel-Wanderung}

% cinnabar color and \stempel command are defined in lernpfad.tex (loaded earlier).

\vspace{-5pt}
{\color{maincolor}\rule{\textwidth}{0.4pt}}
\vspace{4pt}

{\color{maincolor}\textbf{Worum es geht:}}
\begin{itemize}[leftmargin=*]
  \item Auf japanischen Wanderwegen gibt es an jeder Station einen Stempel,
    den du in dein Heft dr\"uckst. Wer alle Stempel hat, hat den Berg geschafft.
  \item Wir machen das mit Code: jede gemeisterte Station gibt
    \textbf{einen Stempel}. Mit jedem Stempel wird ein Teil deines
    \textit{paint-by-numbers}-Bildes sichtbar.
  \item Es gibt insgesamt \textbf{10 Stationen}, vom ersten geladenen
    Pixel bis zur fertigen Abgabe.
\end{itemize}
\vspace{-2pt}
{\color{maincolor}\rule{\textwidth}{0.4pt}}
\vspace{8pt}

% ============================================================================
\section{So funktioniert es}

\begin{enumerate}
  \item Zu Beginn des Projekts bekommst du eine \textbf{Login-Karte} mit
    deinem Benutzernamen und einem Passwort.
  \item Damit \"offnest du die Website \textit{Bildfilter Stempel-Wanderung}
    im Browser und meldest dich an.
  \item Wenn du eine Station geschafft hast, zeigst du dein Resultat der
    Lehrperson.
  \item Die Lehrperson setzt dir den Stempel direkt in dein Heft. Die
    n\"achsten zwei Farben auf deinem Bild werden sofort sichtbar.
\end{enumerate}

\begin{definitionbox}{Login-Karte}
Auf deiner Login-Karte stehen:
\begin{itemize}
  \item dein \textbf{Benutzername} (in der Regel dein Vorname, klein geschrieben),
  \item dein \textbf{Passwort} (sechs Zeichen),
  \item die \textbf{Adresse} der Website.
\end{itemize}
Bewahre die Karte gut auf. Wenn du sie verlierst, sag der Lehrperson Bescheid,
dann gibt es ein neues Passwort.
\end{definitionbox}

% ============================================================================
\section{Der Trail}

\begin{center}
\begin{tikzpicture}[
  station/.style={
    circle, draw=cinnabar, fill=cinnabar!8,
    minimum size=10mm, font=\small\bfseries\color{cinnabar},
    line width=0.9pt, inner sep=0pt,
  },
  trail/.style={cinnabar!55, line width=1.4pt, rounded corners=4pt},
  caption/.style={font=\scriptsize, align=center, color=black!75, text width=24mm},
]
  % Row 1 (left to right)
  \node[station] (s1) at (0, 0)  {1};
  \node[station] (s2) at (3, 0)  {2};
  \node[station] (s3) at (6, 0)  {3};
  \node[station] (s4) at (9, 0)  {4};
  % Row 2 (right to left)
  \node[station] (s5) at (9, -2.5) {5};
  \node[station] (s6) at (6, -2.5) {6};
  \node[station] (s7) at (3, -2.5) {7};
  % Row 3 (left to right)
  \node[station] (s8) at (3, -5)   {8};
  \node[station] (s9) at (6, -5)   {9};
  \node[station] (s10) at (9, -5)  {10};

  % Trail (winding)
  \draw[trail] (s1) -- (s2) -- (s3) -- (s4)
               -- (s5) -- (s6) -- (s7)
               -- (s8) -- (s9) -- (s10);

  % Captions
  \node[caption] at (0, 0.85)   {Erstes\\Bild};
  \node[caption] at (3, 0.85)   {Kein Rot};
  \node[caption] at (6, 0.85)   {Spiegelwelt};
  \node[caption] at (9, 0.85)   {Graustufen};
  \node[caption] at (9, -1.65)  {Heller,\\dunkler};
  \node[caption] at (6, -1.65)  {Schwarz-\\Weiss};
  \node[caption] at (3, -1.65)  {Nachbarn\\mischen};
  \node[caption] at (3, -4.15)  {Spiegelei};
  \node[caption] at (6, -4.15)  {Eigener\\Filter};
  \node[caption] at (9, -4.15)  {Auf dem\\Gipfel};

  % Decorative summit marker after last station
  \node[font=\Large\color{cinnabar}] at (10.4, -5) {$\blacktriangle$};
\end{tikzpicture}
\end{center}

\vspace{6pt}

% ============================================================================
\section{Stationen im Detail}

\renewcommand{\arraystretch}{1.2}

\begin{center}
\begin{tabularx}{\textwidth}{>{\bfseries\color{cinnabar}}c >{\bfseries}p{3.6cm} X p{1.8cm}}
\toprule
\tableheadercolor
\# & Stempel & Anforderung & Lernpfad \\
\midrule
1  & Erstes Bild
   & Bild mit \texttt{lade\_bild} in eine Liste laden, H\"ohe und Breite ausgeben.
   & Schritt 0 \\
2  & Kein Rot
   & Rot-Wert aller Pixel auf $0$ setzen und speichern.
   & Schritt 1 \\
3  & Spiegelwelt
   & \texttt{invertieren(bild)}: jeder Wert wird zu $255 -$ Wert.
   & Schritt 1 \\
4  & Graustufen
   & \texttt{graustufen(bild)}: Mittelwert $\frac{R+G+B}{3}$.
   & Schritt 2 \\
5  & Heller, dunkler
   & \texttt{helligkeit(bild, h)} mit Clamping auf $[0, 255]$.
   & Schritt 2 \\
6  & Schwarz-Weiss
   & \texttt{schwellwert(bild, t)}: Pixel \"uber $t$ werden weiss.
   & Schritt 3 \\
7  & Nachbarn mischen
   & \texttt{box\_blur(bild)}: $3\times 3$-Nachbarschaft mitteln.
   & Schritt 4 \\
8  & Spiegelei
   & \texttt{spiegeln(bild)}: horizontal spiegeln.
   & Schritt 5 \\
9  & Eigener Filter
   & Ein selbst erfundener Filter, der sichtbar funktioniert.
   & Schritt 6 \\
10 & Auf dem Gipfel
   & Code, drei Vorher/Nachher-Bildpaare und Bericht (2 bis 3 A4-Seiten) abgegeben.
   & Schritt 7 \\
\bottomrule
\end{tabularx}
\end{center}

% ============================================================================
\section{Das Bild}

Dein Stempelheft zeigt eine \textit{paint-by-numbers}-Vorlage des Berges
Fuji mit einer japanischen Pagode. Am Anfang ist nur die Vorlage mit den
Zahlen zu sehen. Mit jedem Stempel wird ein St\"uck der Vorlage durch das
\textit{fertige Bild} ersetzt, das, was du am Ende der Wanderung in
voller Farbe siehst.

Insgesamt ist das Bild in zehn Bereiche aufgeteilt, einer pro Stempel.

\closechapter
